% GENERATED FILE. Edit canonical lessons, then run npm run generate:books.
% canonical-source-hash: fnv1a64:85d1c3f886b27f37
% canonical-lessons: BN-C279-jvola, BN-C279-shikha, BN-C279-sphulinggo, BN-C279-pure-jaoya, BN-C279-bisphoron

\chapter{Fire!}
\label{ch:bn-fire}
\hlchaptermodality{279}

\begin{chapteropening}
\textbf{By the end of this chapter:} \emph{I can describe a fire: flames, sparks and an explosion.}

Complete the last lesson of chapter 279: I can describe a fire: flames, sparks and an explosion.
\end{chapteropening}

\section[\texorpdfstring{\bn{জ্বলা} (jvôlā)}{জ্বলা (jvôlā)}]{\bn{জ্বলা} (jvôlā) — to burn; to be alight}

\label{lesson:BN-C279-jvola}



\begin{quote}
Before the new one: say the Bengali for to overturn, then the Bengali for injured.
\end{quote}

\subsection*{You'll want to know: \bn{জ্বলা}}
\textbf{\bn{জ্বলা}} — \emph{jvôlā} — \textquotedblleft{}to burn; to be alight\textquotedblright{}.

\begin{quote}
\bn{যদি} \bn{রাতে} \bn{বাড়িতে} \bn{আলো} \bn{জ্বলে}, \bn{তাহলে} \bn{বুঝব} \bn{যে} \bn{কেউ} \bn{আছে।}
\end{quote}

\begin{quote}
\emph{jôdi rāte bāṛite ālo jvôle, tāhôle bujhbô je keu āchhe.}
\end{quote}

\begin{quote}
If there's a light on in the house at night, I'll know someone is there.
\end{quote}

Past (he/she …): \bn{সেই} \bn{বাতি} \bn{জ্বলল}, \emph{sei bāti jvôllô}.

\begin{cousinweb}
\textbf{\bn{জ্বলা}} (\emph{jvôlā}) comes down from Sanskrit \emph{jvalati}, 'burns brightly' (CDIAL 5306, which lists Bengali \emph{jalā}, 'to burn'). It is the partner of \textbf{\bn{জ্বালানো}}: you light a lamp, and it burns.
\end{cousinweb}

\begin{grammarlens}[title={What you've built}]
One more word for telling what happened.
\end{grammarlens}

\subsection*{Your turn}
\begin{itemize}
\raggedright
  \item \emph{Say it:} \emph{jvôlā}
  \item \emph{Say it:} \emph{jvôlā}, once more
  \item \emph{Say it:} \emph{jôdi rāte bāṛite ālo jvôle, tāhôle bujhbô je keu āchhe.} — If there's a light on in the house at night, I'll know someone is there.
  \item \emph{Recall:} say the Bengali for to overturn, then the Bengali for injured, then say \emph{jvôlā} again
\end{itemize}

\begin{tcolorbox}[breakable,colback=teal!4,colframe=teal!35!black,title={Before you move on}]
What does \bn{জ্বলা} mean? (\textquotedblleft{}To burn; to be alight\textquotedblright{}.) Say it once more.
\end{tcolorbox}
\section[\texorpdfstring{\bn{শিখা} (shikhā)}{শিখা (shikhā)}]{\bn{শিখা} (shikhā) — a flame}

\label{lesson:BN-C279-shikha}



\begin{quote}
Before the new one: say the Bengali for injured, then the Bengali for to be alight.
\end{quote}

\subsection*{You'll want to know: \bn{শিখা}}
\textbf{\bn{শিখা}} — \emph{shikhā} — \textquotedblleft{}a flame\textquotedblright{}.

\begin{quote}
\bn{মোমবাতির} \bn{শিখা} \bn{বাতাসে} \bn{কাঁপছিল।}
\end{quote}

\begin{quote}
\emph{mombātir shikhā bātāse kā̃pchhilô.}
\end{quote}

\begin{quote}
The candle flame was flickering in the breeze.
\end{quote}

\begin{cousinweb}
\textbf{\bn{শিখা}} (\emph{shikhā}) is Sanskrit \emph{śikhā}, a tuft of hair or a crest (Monier-Williams); the pointed top of a fire is its crest, its flame.
\end{cousinweb}

\begin{grammarlens}[title={What you've built}]
One more word for telling what happened.
\end{grammarlens}

\subsection*{Your turn}
\begin{itemize}
\raggedright
  \item \emph{Say it:} \emph{shikhā}
  \item \emph{Say it:} \emph{shikhā}, once more
  \item \emph{Say it:} \emph{mombātir shikhā bātāse kā̃pchhilô.} — The candle flame was flickering in the breeze.
  \item \emph{Recall:} say the Bengali for injured, then the Bengali for to be alight, then say \emph{shikhā} again
\end{itemize}

\begin{tcolorbox}[breakable,colback=teal!4,colframe=teal!35!black,title={Before you move on}]
What does \bn{শিখা} mean? (\textquotedblleft{}A flame\textquotedblright{}.) Say it once more.
\end{tcolorbox}
\section[\texorpdfstring{\bn{স্ফুলিঙ্গ} (sphulinggô)}{স্ফুলিঙ্গ (sphulinggô)}]{\bn{স্ফুলিঙ্গ} (sphulinggô) — a spark}

\label{lesson:BN-C279-sphulinggo}



\begin{quote}
Before the new one: say the Bengali for to be alight, then the Bengali for a flame.
\end{quote}

\subsection*{You'll want to know: \bn{স্ফুলিঙ্গ}}
\textbf{\bn{স্ফুলিঙ্গ}} — \emph{sphulinggô} — \textquotedblleft{}a spark\textquotedblright{}.

\begin{quote}
\bn{একটা} \bn{ছোট} \bn{স্ফুলিঙ্গ} \bn{থেকে} \bn{বাজারে} \bn{আগুন} \bn{লাগল।}
\end{quote}

\begin{quote}
\emph{ekṭā chhoṭo sphulinggô theke bājāre āgun lāglô.}
\end{quote}

\begin{quote}
The market caught fire from one small spark.
\end{quote}

\begin{cousinweb}
\textbf{\bn{স্ফুলিঙ্গ}} (\emph{sphulinggô}) is Sanskrit \emph{sphuliṅga}, 'a spark of fire' (Monier-Williams), taken over whole. It is a written word; in talk a spark is often \bn{ফুলকি}.
\end{cousinweb}

\begin{grammarlens}[title={What you've built}]
One more word for telling what happened.
\end{grammarlens}

\subsection*{Your turn}
\begin{itemize}
\raggedright
  \item \emph{Say it:} \emph{sphulinggô}
  \item \emph{Say it:} \emph{sphulinggô}, once more
  \item \emph{Say it:} \emph{ekṭā chhoṭo sphulinggô theke bājāre āgun lāglô.} — The market caught fire from one small spark.
  \item \emph{Recall:} say the Bengali for to be alight, then the Bengali for a flame, then say \emph{sphulinggô} again
\end{itemize}

\begin{tcolorbox}[breakable,colback=teal!4,colframe=teal!35!black,title={Before you move on}]
What does \bn{স্ফুলিঙ্গ} mean? (\textquotedblleft{}A spark\textquotedblright{}.) Say it once more.
\end{tcolorbox}
\section[\texorpdfstring{\bn{পুড়ে} \bn{যাওয়া} (puṛe jāoyā)}{পুড়ে যাওয়া (puṛe jāoyā)}]{\bn{পুড়ে} \bn{যাওয়া} (puṛe jāoyā) — to burn down; to get burnt}

\label{lesson:BN-C279-pure-jaoya}



\begin{quote}
Before the new one: say the Bengali for a flame, then the Bengali for a spark.
\end{quote}

\subsection*{You'll want to know: \bn{পুড়ে} \bn{যাওয়া}}
\textbf{\bn{পুড়ে} \bn{যাওয়া}} — \emph{puṛe jāoyā} — \textquotedblleft{}to burn down; to get burnt\textquotedblright{}.

\begin{quote}
\bn{আগুনে} \bn{দোকানটা} \bn{পুড়ে} \bn{গেল।}
\end{quote}

\begin{quote}
\emph{āgune dokānṭā puṛe gelô.}
\end{quote}

\begin{quote}
The shop burned down in the fire.
\end{quote}

Past (he/she …): \bn{সেই} \bn{দোকান} \bn{পুড়ে} \bn{গেল}, \emph{sei dokān puṛe gelô}.

\begin{cousinweb}
\textbf{\bn{পুড়ে} \bn{যাওয়া}} (\emph{puṛe jāoyā}) is the -e form of \bn{পোড়া}, to be burnt — CDIAL 8397 lists Bengali \emph{poṛā} and \emph{puṛā}, 'to burn, be burnt' — before \textbf{\bn{যাওয়া}}, which finishes the burning off. \textbf{\bn{পোড়ানো}} is to burn something.
\end{cousinweb}

\begin{grammarlens}[title={What you've built}]
One more word for telling what happened.
\end{grammarlens}

\subsection*{Your turn}
\begin{itemize}
\raggedright
  \item \emph{Say it:} \emph{puṛe jāoyā}
  \item \emph{Say it:} \emph{puṛe jāoyā}, once more
  \item \emph{Say it:} \emph{āgune dokānṭā puṛe gelô.} — The shop burned down in the fire.
  \item \emph{Recall:} say the Bengali for a flame, then the Bengali for a spark, then say \emph{puṛe jāoyā} again
\end{itemize}

\begin{tcolorbox}[breakable,colback=teal!4,colframe=teal!35!black,title={Before you move on}]
What does \bn{পুড়ে} \bn{যাওয়া} mean? (\textquotedblleft{}To burn down; to get burnt\textquotedblright{}.) Say it once more.
\end{tcolorbox}
\section[\texorpdfstring{\bn{বিস্ফোরণ} (bisphorôn)}{বিস্ফোরণ (bisphorôn)}]{\bn{বিস্ফোরণ} (bisphorôn) — an explosion}

\label{lesson:BN-C279-bisphoron}



\begin{quote}
Before the new one: say the Bengali for a spark, then the Bengali for to burn down.
\end{quote}

\subsection*{You'll want to know: \bn{বিস্ফোরণ}}
\textbf{\bn{বিস্ফোরণ}} — \emph{bisphorôn} — \textquotedblleft{}an explosion\textquotedblright{}.

\begin{quote}
\bn{রাতে} \bn{কারখানায়} \bn{একটা} \bn{বিস্ফোরণ} \bn{হল।}
\end{quote}

\begin{quote}
\emph{rāte kārkhānāy ekṭā bisphorôn hôlô.}
\end{quote}

\begin{quote}
There was an explosion at the factory last night.
\end{quote}

\begin{cousinweb}
\textbf{\bn{বিস্ফোরণ}} (\emph{bisphorôn}) is a learned word built from Sanskrit \emph{vi-}, 'apart', and \emph{sphoraṇa}, 'throbbing, quivering' (Monier-Williams); Sanskrit itself has the near word \emph{visphoṭa}, 'a crash'. \bn{বিস্ফোরণ} \bn{হওয়া}, an explosion happens.
\end{cousinweb}

\begin{grammarlens}[title={What you've built}]
One more word for telling what happened.
\end{grammarlens}

\subsection*{Your turn}
\begin{itemize}
\raggedright
  \item \emph{Say it:} \emph{bisphorôn}
  \item \emph{Say it:} \emph{bisphorôn}, once more
  \item \emph{Say it:} \emph{rāte kārkhānāy ekṭā bisphorôn hôlô.} — There was an explosion at the factory last night.
  \item \emph{Recall:} say the Bengali for a spark, then the Bengali for to burn down, then say \emph{bisphorôn} again
\end{itemize}

\begin{tcolorbox}[breakable,colback=teal!4,colframe=teal!35!black,title={Before you move on}]
What does \bn{বিস্ফোরণ} mean? (\textquotedblleft{}An explosion\textquotedblright{}.) Say it once more.
\end{tcolorbox}
